%!TeX root=../thesis.tex

% Keywords must be defined before the abstract/resumo commands.
\keywords{Software Architecture, Architectural Uncertainty, Hypotheses Engineering, ArchHypo, Hypo-Stage, Empirical Evaluation}

\palavraschave{Arquitetura de Software, Incerteza Arquitetural, Engenharia de Hipóteses, ArchHypo, Hypo-Stage, Avaliação Empírica}

% English abstract (primary)
\abstract{
% TODO: revise after chapters are more mature.
Software engineering teams frequently face uncertainties related to their systems'
architecture yet lack adequate methods and tools to manage them explicitly and
collaboratively. The ArchHypo technique proposes using hypotheses engineering to
make such uncertainties visible and manageable; however, its manual application
imposes a significant learning curve. This work empirically evaluates and refines
Hypo-Stage, a support tool for the ArchHypo technique, by applying it with multiple
software engineering teams in a real-world organization with a scaled software
product. The research adopts a mixed-methods design combining analysis of artifacts
produced in the tool, questionnaires, and interviews with team members. Preliminary
results indicate how the tool supports the management of architectural uncertainties
and inform its ongoing improvement.
}

% Portuguese abstract (required by IME-USP regulations)
\resumo{
% TODO: revise after chapters are more mature.
Equipes de engenharia de software frequentemente enfrentam incertezas relacionadas
à arquitetura de seus sistemas, mas carecem de métodos e ferramentas adequadas para
gerenciá-las de forma explícita e colaborativa. A técnica ArchHypo propõe o uso de
engenharia de hipóteses para tornar essas incertezas visíveis e gerenciáveis, porém
sua aplicação manual impõe uma curva de aprendizado significativa. Este trabalho
avalia e refina empiricamente o Hypo-Stage, uma ferramenta de apoio à técnica
ArchHypo, por meio de sua aplicação com múltiplas equipes de engenharia de
software em uma organização real com um produto em escala. A pesquisa adota
um design de métodos mistos, combinando análise de artefatos produzidos na
ferramenta, questionários e entrevistas com membros das equipes. Os resultados
preliminares indicam como a ferramenta apoia o gerenciamento de incertezas
arquiteturais e informam seu aprimoramento contínuo.
}
